\section{Consequences and discussion}\label{sec:discussion}

We close by drawing out what the three theorems say together, and by comparing the techniques with the two-color theory they
extend.

\subsection{The cost of going online}\label{sec:disc-cost}

\Cref{thm:offline,thm:online} let us read off the price the online constraint exacts as a function of the palette size. Define
the \emph{online overhead} at palette $k$ as the ratio of the online value to the offline value,
\[
  \kappa_k\ :=\ \frac{\ondisc_k}{\offdisc_k}.
\]
For fixed $k$, \Cref{thm:separation} gives $\offdisc_k=\Oh(\sqrt d)$ and $\ondisc_k=\Th(\loglog n)$, so
$\kappa_k=\Om(\loglog n/\sqrt d)$, which is $\omega(1)$ across the whole sparsity range $2\le d\le\Hscale^2/\log\Hscale$ and
grows as $d$ shrinks: the online penalty is largest exactly when the instance is sparsest. For growing $k$ the offline floor
falls to $\Th(\Psifn(\lam,\log k))$ while the online value stays at most $\loglog n$, so the upper bracket gives
$\kappa_k\le\Oh(\loglog n/\Psifn(\lam,\log k))$. The qualitative conclusion is robust: more colors are a genuine resource
offline and an illusory one online, and the gap between the two behaviours is the content of the separation.

\begin{corollary}[Offline monotonicity]\label{cor:offline-mono}
For $C_0\log^5(dk)\le k\le d$ with $d,k=(\log n)^{o(1)}$ and $\lam\ge1$, $\offdisc_k=\Th(\Psifn(\lam,\log k))$ is
nonincreasing in $k$ up to constants, and reaches its smallest value $\Th(\log k/\loglog k)$ at the bounded-occupancy end
$\lam=\Th(1)$, i.e.\ $k=\Th(d)$.
\end{corollary}

\begin{proof}
$\Psifn(\lam,\log k)$ with $\lam=\Th(d/k)$ is nonincreasing in $k$ in each of its three regimes of \eqref{eq:Psi-regimes}: the
Gaussian term $\sqrt{\lam\log k}=\Th(\sqrt{(d/k)\log k})$ decreases once $k\ge\log k$, and the two large-deviation terms are
bounded by $\log k/\loglog k$. At $k=\Th(d)$, $\lam=\Th(1)$ and the bounded-occupancy term $\log k/\loglog k$ governs, by
\Cref{thm:offline}.
\end{proof}

At constant $k$ the offline value is the larger $\Th(\sqrt d)$ of the classical regime (\Cref{thm:separation}), so across the
whole palette the offline discrepancy decreases from $\Th(\sqrt d)$ down to the balls-into-bins value
$\Th(\log k/\loglog k)$. The bottoming-out point is the occupancy regime $\lam=\Th(1)$, where each color sees $\Th(1)$
arrivals per coordinate and the discrepancy is the maximum occupancy $\Th(\log d/\loglog d)$, the same law that governs the
longest hash-probe sequence \cite{Gonnet1981}. Within the studied range $k\le d$ this is the smallest the offline value
gets.

\subsection{Why the recursion loses no logarithm}\label{sec:disc-technique}

The single most delicate point in the offline upper bound is that the color-tree composes $\log_2 k$ levels of partial
coloring without paying a factor of $\log k$ (or of $\log d$, or of $\loglog n$) for the composition. Two design choices make
this possible, and both are worth isolating because they are where a naive attempt fails.

First, the splitting error at a node is \emph{density-sensitive}: a node controlling $m$ columns contributes
$\Psifn(\mu,\ell)$ at the occupancy scale of that node, and not a uniform $\sqrt m$. A flat Lovett--Meka rounding at each
node, with a common threshold, would contribute $\sqrt m$ per level and sum to $\Th(\sqrt T\log k)$, off by a large
polynomial factor. The local tuning of thresholds is what keeps each node's error at the right scale.

Second, the errors \emph{attenuate} down the tree. An error created at a node with $q$ leaves below it is divided among those
leaves and reaches any single leaf reduced by $\Oh(1/q)$. The leaf therefore sees a geometric sum $\sum_q \Psifn(\Th(\lam
q),\ell_q)/q$ along its root-to-leaf path, and this sum is dominated by its smallest scales and totals $\Th(\Psifn(\lam,\log
k))$. A union bound over levels would instead have multiplied the per-level error by the number of levels; the attenuation is
exactly what replaces that product by a convergent series. The same two features (a density-sensitive per-step error and a
geometric attenuation) are what one should look for in any recursive rounding that aims to be tight at every scale at once,
and they are the reusable methodological output of the construction.

\subsection{What the online transplant does and does not use}\label{sec:disc-online}

The online upper bound is deliberately a thin layer over \cite{AT2025}. The point of isolating properties (i) and (ii) of the
repair rule (\Cref{sec:online}) is that they are the \emph{entire} interface between the algorithm and the
Altschuler--Tikhomirov analysis: any online rule satisfying them, with any number of colors, inherits the
$\Oh(\loglog n)$ bound. This is what makes the multicolor upper bound short, and it also explains why the online value cannot
drop below the two-color barrier: the lower bound projects $k$ colors onto two by a zero-sum signing, and the projection is a
legitimate online two-coloring whose discrepancy is at most $(k/2)\ondisc_k$. The barrier is therefore not an artifact of the
algorithm but a feature of the model: the irrevocable commitment is what freezes $\Th(\loglog n)$, and no palette unfreezes
it.

The growing-$k$ frontier (\Cref{sec:open}) is where this thin layer stops sufficing. There the question is no longer whether
a coordinate is exceptional but how its $k$ colors are distributed within the exceptional band, a genuinely $k$-dependent
statistic that the two-color analysis does not track. The obstruction we record is precise: a permutation-invariant potential
in $(\rangeop,\maxmult)$ overcounts by a $\log k$ factor, and the asymmetric rule that would avoid the overcount has an
expectation-level tail that the phase conveyor refutes. Pinning the growing-$k$ order is, in our reading, the natural next
target, and the bracket $\max\{\Psifn(\lam,\log k),\loglog n/k\}\le\ondisc_k\le\loglog n$ says exactly how much room is left.
